% !PRINTED 364
\weberpage{364}
%<<CONT
to present. We confine ourselves here to what is required for the attainment of our chief aim, namely the determination of the \gesperrt{Galois} group of the problem.

We start from an arbitrarily selected \gesperrt{Steiner} complex:
\weq{1}{x_1 y_1,\ x_2 y_2,\ x_3 y_3,\ x_4 y_4,\ x_5 y_5,\ x_6 y_6.}
The form of the equation \S.\,96, (5) then shows that in the complex determined by the pair $x_1 x_2$ the pair $y_1 y_2$, and in the complex determined by $x_1 y_2$ the pair $x_2 y_1$, must occur.

We therefore consider, besides (1), the two further complexes:
\weq{2}{x_1 x_2,\quad y_1 y_2,\quad \ldots}
\weq{3}{x_1 y_2,\quad x_2 y_1,\quad \ldots}
Each of the two complexes (2), (3) contains, besides those already known, eight further double tangents, and these must all be different from the $x_i, y_i$, because, if e.g. $x_3$ occurred in (2), then $x_1 x_2 x_3$ would be syzygetic, which contradicts theorem 3., (\S.\,96). Likewise the two complexes (2) and (3) can contain no common double tangent besides $x_1, x_2, y_1, y_2$. For if, say, $z$ occurred in both, then $x_1 x_2 z$ would be syzygetic by (2) and azygetic by (3), which is a contradiction. From this it follows:

\gesperrt{4. In the three complexes (1), (2), (3) taken together all 28 double tangents occur.}

From this it follows that every double tangent which forms a syzygetic triple with any pair $x_1 y_1$ must occur in the complex $x_1 y_1$; that therefore every syzygetic triple is completed by a determinate further double tangent into a syzygetic quadruple; and that consequently a conic which passes through the points of contact of three double tangents cuts the curve $f$ in the points of contact of a fourth double tangent.

Two pairs of a complex always form a syzygetic quadruple, and however one may divide such a quadruple into two pairs, both pairs always belong to the same complex.

Since from 28 things one can form 14.27 pairs, since every pair of double tangents determines a complex, and in every
%CONT>>
